% ============================================================
%  Artificial Intelligence: A Textbook
%  Chapter 4 — Propositional Logic
%  Offline Beamer slides (uses shared theme.tex + branding.tex)
%  Prince Sattam Bin Abdulaziz University — CCES
% ============================================================
\documentclass[aspectratio=169,11pt]{beamer}

\input{theme.tex}

% ---------- COURSE RE-BRANDING (AI instead of CS3401) --------
\graphicspath{{../chapters/images/ch4/}{../chapters/images/}}
% ============================================================
%  Artificial Intelligence: A Textbook
%  Shared re-branding for the AI course (input AFTER theme.tex)
%  Overrides the CS3401 footer/title page from theme.tex.
% ============================================================

% Footer rebranded for the AI course
\setbeamertemplate{footline}{%
  \begin{beamercolorbox}[wd=\paperwidth,ht=0.4pt,dp=0pt]{footer rule}{}
  \end{beamercolorbox}
  \leavevmode
  \hbox{%
    \begin{beamercolorbox}[wd=0.60\paperwidth,ht=1.8em,dp=0.6em,leftskip=1em]{normal text}%
      \textcolor{PrimaryDark}{\tiny\textbf{Dr. Khalil Chebil} %
        \quad\textbar\quad CCES, PSAU}%
    \end{beamercolorbox}%
    \begin{beamercolorbox}[wd=0.40\paperwidth,ht=1.8em,dp=0.6em,rightskip=1em,right]{normal text}%
      \textcolor{PrimaryDark}{\tiny\insertshorttitle%
        \quad\textbar\quad\insertframenumber\,/\,\inserttotalframenumber}%
    \end{beamercolorbox}%
  }%
}

% Title page rebranded for the AI course
\setbeamertemplate{title page}{%
  \begin{tikzpicture}[remember picture,overlay]
    \fill[PrimaryDark](current page.north west) rectangle (current page.south east);
    \fill[PrimaryMid,opacity=0.45]
      ([shift={(2.5cm,-2.5cm)}]current page.north east) circle(5.5cm);
    \fill[AccentOrange,opacity=0.10]
      ([shift={(-1.5cm,2.0cm)}]current page.south west) circle(4.2cm);
    \draw[white,opacity=0.08,line width=40pt]
      ([shift={(0cm,-0.5cm)}]current page.north west)
      to[out=-30,in=210]
      ([shift={(0cm,0.5cm)}]current page.south east);
  \end{tikzpicture}%
  \vspace{0.6cm}
  \begin{center}
    {\tiny\textcolor{PrimaryLight}{\textbf{CCES} \textendash{}
      College of Computer Engineering and Sciences}}\\[0.2em]
    {\tiny\textcolor{PrimaryLight}{Prince Sattam Bin Abdulaziz University}}\\[1.2em]
    {\footnotesize\textcolor{AccentYellow}{\textbf{CS600: Advanced AI }}}\\[1.5em]
    \textcolor{PrimaryMid}{\rule{0.55\textwidth}{0.6pt}}\\[0.9em]
    {\Large\usebeamerfont{title}\textcolor{white}{\inserttitle}}\\[0.4em]
    {\small\usebeamerfont{subtitle}\textcolor{PrimaryLight}{\insertsubtitle}}\\[0.7em]
    \textcolor{PrimaryMid}{\rule{0.55\textwidth}{0.6pt}}\\[1.0em]
    {\small\textcolor{white}{\insertauthor}}\\[0.3em]
    {\footnotesize\textcolor{PrimaryLight}{\insertdate}}
  \end{center}%
}


% ---------- SLIDE METADATA ----------------------------------
\title[Ch.\,4 — Propositional Logic]{Propositional Logic}
\subtitle{Chapter 4 — Syntax, Semantics, and Automated Inference}
\author{Dr.\ Khalil Chebil}
\date{Artificial Intelligence Course}

% Inline check / cross marks
\newcommand{\cmark}{\textcolor{SuccessGreen}{\ensuremath{\checkmark}}}
\newcommand{\xmark}{\textcolor{AlertRed}{\ensuremath{\times}}}
% Truth-value shorthands
\newcommand{\TT}{\textcolor{SuccessGreen}{\textbf{T}}}
\newcommand{\FF}{\textcolor{AlertRed}{\textbf{F}}}

\begin{document}

% ============================================================
%  TITLE
% ============================================================
{\setbeamertemplate{footline}{}
\begin{frame}[plain,noframenumbering]
  \titlepage
\end{frame}}

% ============================================================
%  OUTLINE
% ============================================================
\begin{frame}{Chapter Outline}{What we will cover}
  \begin{columns}[T]
    \begin{column}{0.5\textwidth}
      \tableofcontents[sections={1-3}]
    \end{column}
    \begin{column}{0.5\textwidth}
      \tableofcontents[sections={4-5}]
    \end{column}
  \end{columns}
  \vfill
  \begin{block}{Learning Objectives}
    \begin{itemize}
      \item Master propositional \keyword{syntax}, \keyword{semantics}, and truth tables
      \item Apply \keyword{equivalences} and convert to \keyword{CNF}
      \item Understand \keyword{entailment}, \keyword{model checking}, and \keyword{resolution}
      \item Use \keyword{DPLL} for SAT and \keyword{chaining} for Horn clauses
    \end{itemize}
  \end{block}
\end{frame}

% ============================================================
\section{Foundations \& Syntax}
% ============================================================

\begin{frame}{Why Logic in AI?}{A formal language for knowledge \& reasoning}
  \begin{columns}[T]
    \begin{column}{0.5\textwidth}
      Propositional (Boolean) logic is the foundation of knowledge representation
      and automated reasoning.
      \begin{itemize}\small
        \item \keyword{Declarative} — state \emph{what} is true
        \item \keyword{Compositional} — build complex from simple
        \item \keyword{Unambiguous} — precise semantics
        \item \keyword{Inference} — derive new knowledge
        \item \keyword{Explainable} — transparent steps
      \end{itemize}
    \end{column}
    \begin{column}{0.5\textwidth}
      \begin{exampleblock}{Applications}
        \small
        \begin{itemize}
          \item Expert systems \& diagnosis
          \item Hardware / software verification
          \item Planning \& scheduling
          \item Semantic web, knowledge graphs
          \item Circuit design \& SAT
        \end{itemize}
      \end{exampleblock}
    \end{column}
  \end{columns}
\end{frame}

\begin{frame}{Syntax: Propositions \& Connectives}{Building well-formed formulas}
  \begin{columns}[T]
    \begin{column}{0.52\textwidth}
      \keyword{Atomic propositions} $P, Q, R, \dots$ each denote a statement that
      is \emph{true} or \emph{false}.
      \vspace{0.3em}
      \small
      \renewcommand{\arraystretch}{1.3}
      \begin{tabularx}{\linewidth}{@{}c l >{\raggedright\arraybackslash}X@{}}
        \toprule
        \rowcolor{PrimaryDark}
        \textcolor{white}{Symbol} & \textcolor{white}{Name} & \textcolor{white}{Meaning} \\
        \midrule
        $\neg P$ & NOT & opposite of $P$ \\
        \rowcolor{PrimaryLight!25}
        $P \land Q$ & AND & both \\
        $P \lor Q$ & OR & either or both \\
        \rowcolor{PrimaryLight!25}
        $P \rightarrow Q$ & IF-THEN & if $P$ then $Q$ \\
        $P \leftrightarrow Q$ & IFF & $P$ exactly when $Q$ \\
        \bottomrule
      \end{tabularx}
    \end{column}
    \begin{column}{0.48\textwidth}
      \begin{block}{Well-formed formulas (WFF)}
        \small
        \begin{itemize}
          \item Every atom is a WFF
          \item If $\alpha$ is a WFF, so is $\neg\alpha$
          \item If $\alpha,\beta$ are WFFs, so are
            $(\alpha \land \beta),(\alpha \lor \beta),$
            $(\alpha \rightarrow \beta),(\alpha \leftrightarrow \beta)$
        \end{itemize}
      \end{block}
      \centering\scriptsize
      $\neg(P \land \neg Q)$ \cmark \qquad $P \land \land Q$ \xmark
    \end{column}
  \end{columns}
\end{frame}

% ============================================================
\section{Semantics \& Truth Tables}
% ============================================================

\begin{frame}{Semantics: Models \& Truth Tables}{A model assigns a truth value to every symbol}
  \begin{center}
  \small
  \renewcommand{\arraystretch}{1.15}
  \setlength{\tabcolsep}{4pt}
  \begin{tabular}{cc|c|c|c|c|c}
    \rowcolor{PrimaryDark}
    \textcolor{white}{$P$} & \textcolor{white}{$Q$} &
    \textcolor{white}{$\neg P$} & \textcolor{white}{$P\land Q$} &
    \textcolor{white}{$P\lor Q$} & \textcolor{white}{$P\rightarrow Q$} &
    \textcolor{white}{$P\leftrightarrow Q$} \\
    \midrule
    \TT & \TT & \FF & \TT & \TT & \TT & \TT \\
    \rowcolor{PrimaryLight!20}
    \TT & \FF & \FF & \FF & \TT & \FF & \FF \\
    \FF & \TT & \TT & \FF & \TT & \TT & \FF \\
    \rowcolor{PrimaryLight!20}
    \FF & \FF & \TT & \FF & \FF & \TT & \TT \\
    \bottomrule
  \end{tabular}
  \end{center}
  \vspace{0.4em}
  \begin{columns}[T]
    \begin{column}{0.5\textwidth}
      \begin{block}{Key insight}
        \small $P \rightarrow Q \;\equiv\; \neg P \lor Q$
      \end{block}
    \end{column}
    \begin{column}{0.5\textwidth}
      \begin{exampleblock}{Key insight}
        \small $P \leftrightarrow Q \;\equiv\; (P\rightarrow Q)\land(Q\rightarrow P)$
      \end{exampleblock}
    \end{column}
  \end{columns}
\end{frame}

\begin{frame}{Logical Equivalences}{Same truth value in all models: $\alpha \equiv \beta$}
  \begin{columns}[T]
    \begin{column}{0.5\textwidth}
      \small
      \begin{block}{Core laws}
        \begin{itemize}\small
          \item \keyword{Double negation}: $\neg\neg P \equiv P$
          \item \keyword{Commutative} / \keyword{Associative}
          \item \keyword{Distributive}: \\
            $P \land (Q \lor R) \equiv (P\land Q)\lor(P\land R)$
        \end{itemize}
      \end{block}
    \end{column}
    \begin{column}{0.5\textwidth}
      \small
      \begin{exampleblock}{De Morgan's laws}
        \begin{align*}
          \neg(P \land Q) &\equiv \neg P \lor \neg Q \\
          \neg(P \lor Q)  &\equiv \neg P \land \neg Q
        \end{align*}
      \end{exampleblock}
      \begin{block}{Eliminations}
        \small
        $P \rightarrow Q \equiv \neg P \lor Q$ \\
        $P \leftrightarrow Q \equiv (P\rightarrow Q)\land(Q\rightarrow P)$
      \end{block}
    \end{column}
  \end{columns}
\end{frame}

\begin{frame}{Satisfiability, Validity \& Entailment}{The vocabulary of reasoning}
  \begin{columns}[T]
    \begin{column}{0.5\textwidth}
      \begin{block}{Three properties}
        \small
        \begin{itemize}
          \item \keyword{Satisfiable}: true in \emph{some} model \\ \scriptsize e.g.\ $P \land Q$
          \item \keyword{Unsatisfiable}: false in \emph{all} \\ \scriptsize e.g.\ $P \land \neg P$
          \item \keyword{Valid} (tautology): true in \emph{all} \\ \scriptsize e.g.\ $P \lor \neg P$
        \end{itemize}
      \end{block}
    \end{column}
    \begin{column}{0.5\textwidth}
      \begin{exampleblock}{Entailment $KB \models \alpha$}
        \small $\alpha$ is true in every model where $KB$ is true.
        \\[0.3em]
        $KB: P\rightarrow Q,\ P \quad\Rightarrow\quad KB \models Q$
      \end{exampleblock}
      \begin{alertblock}{Refutation principle}
        \small $KB \models \alpha \iff (KB \land \neg\alpha)$ is \keyword{unsatisfiable}.
      \end{alertblock}
    \end{column}
  \end{columns}
\end{frame}

\begin{frame}{Conjunctive Normal Form (CNF)}{A conjunction of clauses of literals}
  \begin{columns}[T]
    \begin{column}{0.46\textwidth}
      \begin{block}{Definitions}
        \small
        \begin{itemize}
          \item \keyword{Literal}: $P$ or $\neg P$
          \item \keyword{Clause}: disjunction of literals
          \item \keyword{CNF}: conjunction of clauses
        \end{itemize}
      \end{block}
      \centering\scriptsize
      $(P\lor Q)\land(\neg P\lor R)$ \cmark \\
      $(P\land Q)\lor R$ \xmark
    \end{column}
    \begin{column}{0.54\textwidth}
      \begin{exampleblock}{Conversion algorithm}
        \small
        \begin{enumerate}
          \item Eliminate $\rightarrow,\leftrightarrow$
          \item Push $\neg$ inward (De Morgan)
          \item Distribute $\lor$ over $\land$
        \end{enumerate}
      \end{exampleblock}
      \scriptsize
      \textbf{Example} — convert $\neg(P\rightarrow Q)$:
      \[ \neg(\neg P\lor Q) \;\Rightarrow\; \neg\neg P\land\neg Q \;\Rightarrow\; P\land\neg Q \]
    \end{column}
  \end{columns}
\end{frame}

% ============================================================
\section{Inference: Model Checking, Resolution, DPLL}
% ============================================================

\begin{frame}{Model Checking}{Enumerate all models}
  \begin{columns}[T]
    \begin{column}{0.55\textwidth}
      \keyword{TT-Entails?} checks $KB \models \alpha$ by trying every assignment of
      the $n$ symbols.
      \begin{itemize}\small
        \item For each model: if $KB$ true, require $\alpha$ true
        \item \keyword{Sound} and \keyword{complete}
      \end{itemize}
      \begin{alertblock}{Cost}
        \small $\mathcal{O}(2^n)$ — only practical for small formulas.
      \end{alertblock}
    \end{column}
    \begin{column}{0.45\textwidth}
      \begin{center}
      \begin{tikzpicture}[level distance=0.9cm, level 1/.style={sibling distance=2.4cm},
          level 2/.style={sibling distance=1.2cm},
          n/.style={circle, draw=PrimaryMid, fill=PrimaryLight!45, font=\scriptsize, minimum size=0.5cm},
          l/.style={rectangle, rounded corners=1pt, draw=PrimaryMid, fill=white, font=\scriptsize\ttfamily, inner sep=2pt},
          edge from parent/.style={draw=PrimaryMid}]
        \node[n] {$P$}
          child {node[n] {$Q$}
            child {node[l] {TT}}
            child {node[l] {TF}}
            edge from parent node[left,font=\scriptsize] {T}}
          child {node[n] {$Q$}
            child {node[l] {FT}}
            child {node[l] {FF}}
            edge from parent node[right,font=\scriptsize] {F}};
      \end{tikzpicture}\\
      \scriptsize $2^n$ leaves for $n$ symbols
      \end{center}
    \end{column}
  \end{columns}
\end{frame}

\begin{frame}{Resolution Theorem Proving}{One sound, refutation-complete rule}
  \begin{columns}[T]
    \begin{column}{0.5\textwidth}
      From two clauses with complementary literals, infer their resolvent:
      \[ \frac{\ell \lor A,\quad \neg\ell \lor B}{A \lor B} \]
      \begin{block}{Example rule}
        \small \[ \frac{P \lor Q,\quad \neg P \lor R}{Q \lor R} \]
      \end{block}
      \centering\scriptsize
      Prove $KB \models \alpha$ by showing $KB \land \neg\alpha$ derives the
      \keyword{empty clause} $\square$.
    \end{column}
    \begin{column}{0.5\textwidth}
      \begin{exampleblock}{Refutation: $KB \models R$}
        \small $KB=\{P\rightarrow Q,\ Q\rightarrow R,\ P\}$
      \end{exampleblock}
      \begin{center}
      \begin{tikzpicture}[node distance=0.4cm and 0.5cm,
          c/.style={rectangle, rounded corners=2pt, draw=PrimaryMid, fill=PrimaryLight!40, font=\scriptsize, minimum height=0.5cm},
          r/.style={rectangle, rounded corners=2pt, draw=AccentOrange, fill=AccentYellow!30, font=\scriptsize, minimum height=0.5cm}]
        \node[c] (p) {$P$};
        \node[c, right=1.1cm of p] (pq) {$\neg P\lor Q$};
        \node[r, below=0.5cm of $(p)!0.5!(pq)$] (q) {$Q$};
        \node[c, right=0.5cm of pq] (qr) {$\neg Q\lor R$};
        \node[r, below=0.5cm of $(q)!0.5!(qr)$] (rr) {$R$};
        \node[c, right=0.5cm of qr] (nr) {$\neg R$};
        \node[r, below=0.5cm of $(rr)!0.5!(nr)$, fill=AlertRed!20, draw=AlertRed] (box) {$\square$};
        \draw[-{Latex[length=1.2mm]},PrimaryMid] (p)--(q);
        \draw[-{Latex[length=1.2mm]},PrimaryMid] (pq)--(q);
        \draw[-{Latex[length=1.2mm]},PrimaryMid] (q)--(rr);
        \draw[-{Latex[length=1.2mm]},PrimaryMid] (qr)--(rr);
        \draw[-{Latex[length=1.2mm]},PrimaryMid] (rr)--(box);
        \draw[-{Latex[length=1.2mm]},PrimaryMid] (nr)--(box);
      \end{tikzpicture}\\
      \scriptsize Empty clause $\Rightarrow$ proved.
      \end{center}
    \end{column}
  \end{columns}
\end{frame}

\begin{frame}{DPLL Algorithm for SAT}{The engine behind modern SAT solvers}
  \begin{columns}[T]
    \begin{column}{0.5\textwidth}
      Backtracking search with powerful heuristics:
      \begin{itemize}\small
        \item \keyword{Early termination} — stop when decided
        \item \keyword{Pure symbol} — one polarity $\Rightarrow$ assign it
        \item \keyword{Unit clause} — single literal must be true
        \item \keyword{Splitting} — try both values, backtrack
      \end{itemize}
    \end{column}
    \begin{column}{0.5\textwidth}
      \begin{block}{Properties}
        \small
        \begin{itemize}
          \item Much faster than plain enumeration
          \item Still exponential worst-case
          \item Foundation of \keyword{CDCL} SAT solvers
        \end{itemize}
      \end{block}
      \centering\scriptsize
      Unit propagation + pure literals prune huge parts of the search tree.
    \end{column}
  \end{columns}
\end{frame}

% ============================================================
\section{Horn Clauses \& Chaining}
% ============================================================

\begin{frame}{Horn Clauses}{Restricted but tractable}
  \begin{columns}[T]
    \begin{column}{0.5\textwidth}
      A \keyword{Horn clause} has \emph{at most one} positive literal.
      \begin{itemize}\small
        \item \keyword{Definite}: exactly one positive \\
          \scriptsize $\neg P \lor \neg Q \lor R \equiv (P \land Q \rightarrow R)$
        \item \keyword{Fact}: a single positive literal $P$
        \item \keyword{Goal}: no positive literals
      \end{itemize}
    \end{column}
    \begin{column}{0.5\textwidth}
      \begin{exampleblock}{Why they matter}
        \small
        \begin{itemize}
          \item Many real KBs are naturally Horn
          \item Inference is \keyword{linear time}
          \item Foundation of \keyword{Prolog} / logic programming
        \end{itemize}
      \end{exampleblock}
    \end{column}
  \end{columns}
\end{frame}

\begin{frame}{Forward vs.\ Backward Chaining}{Two ways to reason with Horn clauses}
  \begin{columns}[T]
    \begin{column}{0.5\textwidth}
      \begin{block}{Forward (data-driven)}
        \small Start from facts, fire rules, derive \emph{all} consequences until
        the query appears.
      \end{block}
      \begin{exampleblock}{Backward (goal-driven)}
        \small Start from the query, recurse on rule premises until grounded in facts.
      \end{exampleblock}
    \end{column}
    \begin{column}{0.5\textwidth}
      \small
      \renewcommand{\arraystretch}{1.3}
      \begin{tabularx}{\linewidth}{@{}>{\bfseries}l >{\raggedright\arraybackslash}X >{\raggedright\arraybackslash}X@{}}
        \toprule
        \rowcolor{PrimaryDark}
        \textcolor{white}{Aspect} & \textcolor{white}{Forward} & \textcolor{white}{Backward} \\
        \midrule
        Direction & Data-driven & Goal-driven \\
        \rowcolor{PrimaryLight!25}
        Strategy & Breadth-first & Depth-first \\
        Space & More & Less \\
        \rowcolor{PrimaryLight!25}
        Best for & All facts & Specific query \\
        \bottomrule
      \end{tabularx}
      \centering\scriptsize Both complete for definite clauses; $\mathcal{O}(pn)$.
    \end{column}
  \end{columns}
\end{frame}

% ============================================================
\section{Summary}
% ============================================================

\begin{frame}{Inference Methods Compared}{Choose by problem shape}
  \small
  \renewcommand{\arraystretch}{1.3}
  \begin{center}
  \begin{tabularx}{\textwidth}{@{}>{\bfseries}l cc l >{\raggedright\arraybackslash}X@{}}
    \toprule
    \rowcolor{PrimaryDark}
    \textcolor{white}{Method} & \textcolor{white}{Compl.} & \textcolor{white}{Sound} &
    \textcolor{white}{Complexity} & \textcolor{white}{Best for} \\
    \midrule
    Truth tables & Yes & Yes & $\mathcal{O}(2^n)$ & Small formulas \\
    \rowcolor{PrimaryLight!25}
    Resolution & Yes$^\ast$ & Yes & Exponential & General CNF \\
    DPLL & Yes & Yes & Exponential & SAT problems \\
    \rowcolor{PrimaryLight!25}
    Forward chain & Yes$^{\ast\ast}$ & Yes & $\mathcal{O}(pn)$ & Derive all facts \\
    Backward chain & Yes$^{\ast\ast}$ & Yes & $\mathcal{O}(pn)$ & Specific queries \\
    \bottomrule
  \end{tabularx}
  \end{center}
  \centering\scriptsize $^\ast$refutation-complete \quad $^{\ast\ast}$for Horn clauses
\end{frame}

\begin{frame}{Chapter Summary}{Key takeaways}
  \begin{columns}[T]
    \begin{column}{0.5\textwidth}
      \begin{block}{Represent}
        \small
        \begin{itemize}
          \item Syntax: propositions + connectives
          \item Semantics: models \& truth tables
          \item \keyword{CNF} as the standard inference form
        \end{itemize}
      \end{block}
      \begin{exampleblock}{Reason}
        \small Model checking, \keyword{resolution}, \keyword{DPLL}, and chaining.
      \end{exampleblock}
    \end{column}
    \begin{column}{0.5\textwidth}
      \begin{alertblock}{Limitations}
        \small Propositional logic cannot express \keyword{quantifiers},
        \keyword{objects}, \keyword{relations}, or \keyword{functions} — it must
        enumerate every instance.
      \end{alertblock}
      \begin{block}{Looking ahead}
        \small Ch.\,5: \keyword{First-Order Logic} adds quantifiers and relations.
      \end{block}
    \end{column}
  \end{columns}
\end{frame}

{\setbeamertemplate{footline}{}
\begin{frame}[plain,noframenumbering]
  \begin{tikzpicture}[remember picture,overlay]
    \fill[PrimaryDark](current page.north west) rectangle (current page.south east);
    \fill[PrimaryMid,opacity=0.35]
      ([shift={(0cm,-1.5cm)}]current page.north) circle(6.5cm);
  \end{tikzpicture}
  \vfill
  \begin{center}
    {\LARGE\textcolor{white}{\textbf{Thank you!}}}\\[0.6em]
    {\small\textcolor{PrimaryLight}{Questions \& Discussion}}\\[1.2em]
    {\footnotesize\textcolor{AccentYellow}{Next: Chapter 5 — First-Order Logic}}
  \end{center}
  \vfill
\end{frame}}

\end{document}
